\section{提交材料与公开边界说明}

\subsection{提交材料映射}
初赛提交材料包括作品视频、作品设计报告和重要代码压缩包。根据提交要求，视频控制在 3 分钟以内，报告 PDF 控制在 50 MB 以内，重要代码压缩包控制在 100 MB 以内。项目复杂度远高于单一演示，因此报告采用高密度图表和分章节工程说明，视频则以倍速和多窗口方式展示关键证据。
\begin{figure}[H]
  \centering
  \includegraphics[width=\linewidth]{fig_xrd_submission_boundary_html.png}
  \caption{初赛提交材料与公开边界。}
  \label{fig:xrd_submission}
\end{figure}

\subsection{重要代码包边界}
重要代码包采用“真实项目代码、非核心公开边界”的原则：上传可审查的公网展示站静态前端、报告图表生成脚本、工位碰撞互锁与仿真逻辑、接口 schema、公开演示适配层、公开证据截图、报告渲染页、晶体结构公开缓存和离线演示传感器帧。代码包保留接近提交上限的工程证据体量，同时不公开以下内容：
\begin{itemize}
  \item API Key、账号、Cookie、SSH、WiFi、隧道配置、内网地址和任何可登录凭据；
  \item GGUF、LoRA、BPU bin、tokenizer、tensor 包等模型权重；
  \item 未发表实测数据、原始谱图、failure patterns、私有实验记录和完整训练数据；
  \item 核心预测引擎、可直接复现实验室策略的私有规则库和部署脚本；
  \item 可直接控制底盘、机械臂、升降台、电磁铁或舵机的危险执行脚本；
  \item 任何学校、个人、导师或账号身份信息。
\end{itemize}

\subsection{GitHub 字段}
GitHub 字段填写项目公开边界仓库：
\begin{center}
  \url{https://github.com/Xiaomiju-x/xrd}
\end{center}
该仓库用于提交审查和技术展示，内容限定为真实项目的非核心代码、文档、接口 schema 和可离线查看的前端资产；不上传完整工程、模型权重、私有数据、部署凭据和核心材料预测引擎。报告、视频与仓库对外口径保持一致：强调真机实测、工程验证、shadow mode、安全降级和可追溯证据链。

\begin{warnbox}
报告、视频和公开代码包统一避免绝对化、宣传式或过度承诺措辞，统一使用“真机实测”“工程验证”“本项目实现”“shadow mode”“安全降级”等可防守表述。
\end{warnbox}
